\label{chap:p3-83-45-completitud_finita}

\section{Conservative Extension of Finite Completeness via Projective Multiplicities}

A concrete HMT tower consists of finite sets, transition relations,
restriction maps, and readers. All of this can be
encoded in ZF or ZFC. The collection of all towers constitutes a
definable class; it is unnecessary to postulate that HMT already forms an
autonomous class theory comparable to NBG.

\begin{proposition}[Conservative Coding]
Every finite or countable genealogical object used in this notebook can
be represented by sets and relations in ZFC. Therefore, the
formalizations constructed here do not prove a contradiction in ZFC.
\end{proposition}

\begin{proof}
Each level \(X_m\) is a set, each restriction \(p_{n,m}\) is a
subset of \(X_n\times X_m\), and a countable tower is a function with
domain \(\mathbb N\) encoding these data. Types, readers, and memories
are added as relations or functions. The inverse limit is a subset
of the product \(\prod_mX_m\), which ZFC constructs by Replacement, Union,
Power Set, and Separation.
\end{proof}

What HMT calls into question is not the legitimacy of that encoding, but the practice
of forgetting the maps and retaining only a shadow. If, for example, one applies
the reader
\[
  L:\OmegaInf\longrightarrow\mathbb R
\]
or the cardinal profile
\[
  \mathbb X\longmapsto(|X_m|)_{m\ge0},
\]
nonisomorphic objects may be collapsed. The projection is valid in its
codomain; what would be invalid is to attribute faithfulness to it without proof.

\begin{remark}[Conditions of validity]
It is not claimed that HMT internally satisfies every schema of ZF, IZF, or
NBG. Separation, Replacement, the full Power Set axiom, and a general transfinite
hierarchy have not been verified jointly within an autonomous HMT axiomatics.
\end{remark}

\section{Category of genealogical objects}

\begin{definition}
A genealogical object is
\[
  \mathbb X=
  \bigl((X_m)_{m\ge0},(p_{n,m})_{n\ge m},
  \mathcal T,\mathcal O,\mathcal M\bigr),
\]
where \(\mathcal T\) records types, \(\mathcal O\) observables, and
\(\mathcal M\) memories. A morphism
\(f:\mathbb X\to\mathbb Y\) is a family \(f_m:X_m\to Y_m\) satisfying
\[
  p^Y_{n,m}f_n=f_mp^X_{n,m}
\]
and preserves the indicated data.
\end{definition}

The category thus obtained is a concrete presentation of enriched
profinite objects. The global-sections functor
\[
  \Gamma(\mathbb X)=\varprojlim_mX_m
\]
recovers the complete histories. The small path category of
\(\TPK\) also admits the presheaf envelope
\[
  \widehat{\mathcal P_{\TPK}}
  =
  \mathbf{Set}^{\mathcal P_{\TPK}^{op}}.
\]
This construction provides contextual logic and data sheaves, but does not
warrant declaring Booleanity, choice, or enough points without studying the
specific topos.

\section{Foundations of Prefixes and Coinduction of Sections}

The prefix graph is graded by depth:
\[
  \operatorname{rk}(x_m)=m.
\]
Applying the parent strictly decreases the rank. There are no parental cycles.
\[
  x_m\to x_{m-1}\to\cdots\to x_m
\]
compatible with that graduation.

\begin{proposition}
Well-foundedness of the prefix graph and coinductive existence of an
infinite section are compatible.
\end{proposition}

\begin{proof}
Well-foundedness concerns the direction of descent through parents.
Coinduction concerns a family \((x_m)\) compatible with all descents.
No \(x_m\) belongs to itself or reappears as the same node; the
section is a function on \(\mathbb N\), not a membership cycle.
\end{proof}

This distinction clarifies the return \(9\to1\). The phase may recur, but
the memory coordinate increases. A loop in the phase projection is not a
cycle in the graded state graph.

\section{Fibration of ordinals by compatible genealogical histories}

Von Neumann represents each ordinal as the set of its predecessors,
\citep{vonneumann1923}. HMT may lift that representation without modifying it.

\begin{definition}[Genealogical elevation of a finite ordinal]
Let \(\boldsymbol x=(x_m)\) be a section. Define
\[
  \widehat n_{\boldsymbol x}
  =
  \{(k,x_k):k<n\}.
\]
\end{definition}

\begin{proposition}
The projection onto the first coordinate satisfies
\[
  \pi(\widehat n_{\boldsymbol x})=n,
\]
and the sequence obeys
\[
  \widehat{n+1}_{\boldsymbol x}
  =
  \widehat n_{\boldsymbol x}\cup\{(n,x_n)\}.
\]
\end{proposition}

\begin{proof}
Both identities follow directly from the definition. The added information
does not alter the underlying order; it distinguishes the routes that
realize the same ordinal.
\end{proof}

For \(\omega\) one writes
\[
  \widehat\omega_{\boldsymbol x}
  =
  \{(n,x_n):n\in\mathbb N\}.
\]
The support of each history remains countable, although the space of
all histories may have continuum cardinality when branching
persists.

\begin{remark}[Extension and open domain]
The extension
\[
S_{\alpha+1}=\operatorname{Ext}_\alpha(S_\alpha),
\qquad
S_\lambda=\varprojlim_{\beta<\lambda}S_\beta
\]
is a well-typed transfinite programme, not a result already proved for
every ordinal. The coherence of the limits, the sizes of the
levels, and the existence of extensions must be verified.
\end{remark}

\section{Decategorification Functor and the Cardinality of Published Sections}

Two towers can satisfy
\[
  |X_m|=|Y_m|\qquad\text{for every }m
\]
and fail to be isomorphic: it suffices that their restriction maps have distinct
branching patterns. The cardinal profile is therefore a useful
but nonfaithful invariant.

\begin{remark}[Structural Interpretation]
Counting how many cells exist does not say which cell prolongs which, which
orientation is reversed, which carry is transported, or which state returns. The
cardinal measures the inventory; the genealogy measures the organization.
\end{remark}

\begin{remark}[Separation of results and scope]
The ZFC coding, the section as a limit, and the distinction between observational
equality and state identity are
\emph{results established in the preceding sections}. The unified category, the presheaf
topos, and the fibered ordinal are
\emph{explicit constructions of the present development}. No complete autonomous HMT theory
of classes or transfinite ordinals is proclaimed.
\end{remark}


\section{Profinite limit of the system of genealogical truncations}

Let
\[
  X=\varprojlim_mX_m
\]
an inverse limit of finite discrete sets with surjective maps.
We equip \(X\) with the initial topology of the projections
\(\pi_m:X\to X_m\). The cylinders
\[
  [A]_m=\pi_m^{-1}(A),
  \qquad A\subseteq X_m,
\]
are clopen and form a basis.

\begin{theorem}[Cylindrical power]
\label{p3-83-c45-s2:thm:clopen-power}
\[
  \Clop(X)
  \cong
  \varinjlim_m\bigl(\mathcal P(X_m),p_{n,m}^{-1}\bigr).
\]
\end{theorem}

\begin{proof}
Every cylinder is clopen. Conversely, if \(C\subset X\) is clopen, then for
each \(x\in C\) there is a cylinder contained in \(C\). Compactness of \(C\)
gives a finite subcover. Lifting those cylinders to a common
depth \(m\) yields \(C=\pi_m^{-1}(A)\) for some
\(A\subseteq X_m\). Two representatives at different depths define the
same clopen set exactly when they agree after a common lift.
\end{proof}

\begin{theorem}[Closed-form power]
\label{p3-83-c45-s2:thm:closed-power}
\[
  \Closed(X)
  \cong
  \varprojlim_m\bigl(\mathcal P(X_m),p_{n,m*}\bigr),
\]
where \(p_*(A)=p(A)\).
\end{theorem}

\begin{proof}
If \(F\subset X\) is closed, it is compact. Its projections
\(F_m=\pi_m(F)\) satisfy \(p_{n,m}(F_n)=F_m\). Conversely, a compatible family
\((F_m)\) defines the inverse limit
\[
  F=\{x\in X:\pi_m(x)\in F_m\text{ for every }m\}.
\]
It is an intersection of closed sets and, by the surjective compatibility
of the restrictions \(F_n\to F_m\), its projections are exactly
\(F_m\). The two constructions are inverse.
\end{proof}

\section{Obstruction to the Reconstruction of the Genealogical Power from Finite Reductions}

In a Cantor space,
\[
 |\Clop(X)|=\aleph_0,
 \qquad
 |\Closed(X)|=2^{\aleph_0},
 \qquad
 |\mathcal P(X)|=2^{2^{\aleph_0}}.
\]
The three objects cannot be identified.

This separation does not presuppose the continuum hypothesis. Writing
\(\mathfrak c=2^{\aleph_0}\),
\[
 |\Clop(X)|=\aleph_0,\qquad
 |\Closed(X)|=\mathfrak c,\qquad
 |\mathcal P(X)|=2^{\mathfrak c}.
\]
\(\mathsf{CH}\) identifies only \(\mathfrak c\) with \(\aleph_1\); it does not
collapse any of the three layers. Gödel--Cohen independence changes which
cardinal may occupy \(\mathfrak c\), not Cantor's inequality
\(\mathfrak c<2^{\mathfrak c}\).

\begin{proposition}[Dense falsifier]
If \(D\subsetneq X\) is dense, then
\[
  \pi_m(D)=X_m
  \qquad\text{for every }m,
\]
so that \(D\) and \(X\) possess the same finite shadows.
\end{proposition}

\begin{proof}
Every nonempty cylindrical fiber \(\pi_m^{-1}(a)\) is open. By density, it
intersects \(D\). Hence \(a\in\pi_m(D)\) for every
\(a\in X_m\).
\end{proof}

Finite shadows reconstruct the closure of a subset, not the
arbitrary subset itself. This is the precise point at which an identification
\[
  \varprojlim_m\mathcal P(X_m)
  =
  \mathcal P(\varprojlim_mX_m)
\]
would become false.

\begin{figure}[H]
\centering
\begin{tikzpicture}[>=Latex,node distance=10mm and 17mm,
  box/.style={rounded corners,draw=hmtblue,fill=hmtlightblue,
    minimum width=38mm,minimum height=15mm,align=center},
  full/.style={rounded corners,draw=hmtviolet,fill=white,
    minimum width=40mm,minimum height=15mm,align=center}]
  \node[box] (fin) {\(\mathcal P(X_m)\)\\\scriptsize finite predicates};
  \node[box,right=of fin] (clop) {\(\Clop(X)\)\\\scriptsize direct limit};
  \node[box,below=of clop] (closed) {\(\Closed(X)\)\\\scriptsize inverse limit};
  \node[full,right=of clop] (power) {\(\mathcal P(X)\)\\\scriptsize potencia plena translated};
  \draw[->,thick,hmtblue] (fin)--node[above]{\(p^{-1}\)}(clop);
  \draw[->,thick,hmtblue] (fin) to[bend right=24]
    node[below left]{\(p_*\)} (closed);
  \draw[->,thick,hmtgreen] (clop)--node[above]{\(\subset\)}(power);
  \draw[->,thick,hmtgreen] (closed)--node[below]{\(\subset\)}(power);
  \draw[dashed,very thick,hmtred] (closed.south east)
    .. controls +(8mm,-8mm) and +(-8mm,-8mm) ..
    node[midway,below=2mm]{\(\ne\)} (power.south west);
  \node[below=18mm of closed,text width=.82\textwidth,align=center] {
    \small A proper dense subset and \(X\) have the same finite projections:
    finite history recovers closure, not the full power set.};
\end{tikzpicture}
\caption{Three notions of power that must not be flattened.}
\label{p3-83-c45-s2:fig:three-powers}
\end{figure}

\section{\(4\) choice requirements in the projective system}

Let \(p:X\twoheadrightarrow Y\). It is useful to distinguish:

\begin{enumerate}
\item a \emph{section conjuntista} \(s:Y\to X\) With \(ps=\id_Y\);
\item a \emph{natural family} of sections that commutes with the
restrictions of a tower;
\item an \emph{algebraic section} that also preserves the operations;
\item a \emph{section efectiva} calculada by a algoritmo total.
\end{enumerate}

Every natural, algebraic, or effective section is, upon forgetting structure, a
set-theoretic section. The three additional properties are generally
incomparable: an algorithm may select representatives without being a homomorphism;
a homomorphism may not be computable in the chosen presentation; a
natural family may not preserve the algebraic operation.

\begin{remark}[Conditions of validity]
The axiom of choice guarantees a set-theoretic selection within its
domain. It does not guarantee naturality, algebraicity, effectiveness, or uniform
cost. The independence of \(\mathsf{AC}\) from \(\mathsf{ZF}\)
likewise does not turn these four requirements into one: even in a model
with choice, the structural section requires its own proof.
\end{remark}

\Needspace{36\baselineskip}
\section{Pell Obstruction to the Uniform Periodicity of the Prefix System}

The active generator constructs
\[
  K_\infty\cong C_4\times\mathbb Z_3
\]
and finite projections \(p_m:K_\infty\to K_m\). A
representative of each class may be chosen, but not a homomorphic section compatible with the
group.

\begin{proposition}
There is no homomorphism \(s_m:K_m\to K_\infty\) with
\(p_ms_m=\id\) for the quotient that introduces finite ternary torsion into
\(K_m\).
\end{proposition}

\begin{proof}
A homomorphic section would send an element of order \(3^r\) of \(K_m\) to an
element of the same order in the component \(\mathbb Z_3\). But the additive group
\(\mathbb Z_3\) is torsion-free: if \(3^rx=0\), then \(x=0\).
The image cannot project back onto a nonzero element of order
\(3^r\).
\end{proof}

The obstruction does not say «it cannot be chosen». It says: it cannot be chosen
while simultaneously preserving the projection and the group law. The carry
information that disappears in the quotient reappears as an obstruction to the section.

\begin{remark}[Separation of results and scope]
The profinite continuum and the Pell obstruction are
\emph{results established in the preceding sections}. The two power theorems and the table of
choices are \emph{explicit constructions of the present development}. The dense falsifier precludes
turning an incomplete search for a subset into a power that has already been
reconstructed.
\end{remark}
